\documentclass[11pt]{article}
\usepackage{mathunal-base}
\mathunalfuente{serif}
\providecommand{\nota}[1]{\par\smallskip\noindent{\small\itshape\color{unalblue}#1}\par\smallskip}
\newlist{opciones}{enumerate}{1}
\setlist[opciones]{label=(\roman*),itemsep=1pt,topsep=2pt,leftmargin=1.8em,font=\normalfont}
\newlist{literales}{enumerate}{1}
\setlist[literales]{label=(\alph*),itemsep=1pt,topsep=2pt,leftmargin=1.8em,font=\normalfont}
\newcommand{\rp}{\underline{\hspace{3cm}}}

\renewcommand{\examtitulo}{Ecuaciones Diferenciales (1000007) --- Segundo Examen Parcial}
\renewcommand{\examfecha}{4 de Marzo del 2019}

\begin{document}

\examheader

\vspace{4pt}
\noindent\textit{Duración: 1 hora y 45 minutos. El examen consta de cuatro preguntas en dos hojas impresas por ambos lados; antes de empezar verifique que su examen esté completo y que sus datos sean correctos. No desengrape su examen ni añada otras grapas o su examen será anulado. En las preguntas con procedimiento justifique sus respuestas en los espacios asignados. No está permitido hacer preguntas, usar calculadoras, sacar hojas en blanco ni ningún tipo de apuntes durante el examen. Por favor verifique que su celular esté apagado.}

\vspace{6pt}
\noindent\textbf{1. (20\%)} Dado que $y_1=x^2$ y $y_2=\dfrac{1}{x}$ son soluciones de la ecuación diferencial homogénea asociada a la ecuación diferencial
\[
x^2y''-2y=\frac{5}{x}, \qquad x>0, \tag{*}
\]
usando el método de variación de parámetros halle una solución particular de la ecuación diferencial (*) en $(0,\infty)$.

\vspace{6cm}

\newpage

\noindent\textbf{2. (10\%)} Considere la ecuación diferencial
\[
y''+2by'+b^2y=0,
\]
en la cual $b>0$ es constante y $y(t)$ es la incógnita. Demuestre que todas las soluciones de la ecuación diferencial tienden a cero cuando $t$ tiende a más infinito.

\vspace{6cm}

\noindent\textbf{3. (35\%) Aplicaciones}

\vspace{0.2cm}
\noindent\textbf{a) (15\%)} La posición de un cuerpo, cuya masa es 1 slug, que está sujeto a un resorte, está dada por
\[
x(t)=\frac{1}{3}\cos(2t)-\frac{1}{3}\sin(2t)
\]
pies, para cada instante $t\geq 0$ (medido en segundos).
\begin{opciones}
\item \textbf{(5\%)} Al escribir la posición en la forma $x(t)=A\sin(2t+\Phi)$, los valores de $A$ y $\Phi$ son $A=$~\rp\ y $\Phi=$~\rp
\item \textbf{(5\%)} ¿Cuántos ciclos completos habrá realizado la masa al final de $6\pi$ segundos?
\item \textbf{(5\%)} ¿En qué momento pasa el cuerpo por posición de equilibrio dirigiéndose por segunda vez hacia abajo?
\end{opciones}

\newpage

\noindent\textbf{b) (20\%)} Una masa que pesa 32 lb estira un resorte 8 pies. El sistema masa-resorte se ubica en un medio que presenta una resistencia que es numéricamente igual a 4 veces la velocidad instantánea. En el instante inicial la masa se libera desde un punto situado a $\frac{1}{2}$ pie por encima de la posición de equilibrio, con una velocidad inicial de 4 pies/seg hacia abajo.
\begin{opciones}
\item \textbf{(8\%)} Encuentre la posición $x(t)$ de la masa en cualquier instante $t\geq 0$.
\item \textbf{(4\%)} Determine si la masa pasa por la posición de equilibrio. Si pasa, determine el tiempo o tiempos en que la masa pasa por la posición de equilibrio.
\item \textbf{(4\%)} En el instante $t=3$ seg, determine la posición de la masa e indique hacia dónde se dirige.
\item \textbf{(4\%)} Bosqueje la gráfica del movimiento $t$ contra $x(t)$. Indique claramente el sentido en el cual $x$ es positiva.
\end{opciones}
\textit{(Recuerde que la constante $g$ es 32 pies sobre segundo al cuadrado)}

\vspace{5.5cm}

\newpage

\noindent\textbf{4. (35\%)} En los literales a), b), c) y d) complete según corresponda. Sólo se calificará la respuesta, no es necesario que escriba el procedimiento.

\vspace{0.25cm}
\noindent\textbf{a) (7\%)} Si $r_1=1$ y $r_2=2-2i$ son raíces de la ecuación auxiliar (o característica) asociada a una ecuación diferencial lineal homogénea con coeficientes constantes de orden tres, entonces la ecuación diferencial es: \rp

\vspace{0.3cm}
\noindent\textbf{b) (10\%)} Una forma apropiada de solución particular, según el método de coeficientes indeterminados, de la ecuación diferencial
\[
y^{(6)}+10y^{(4)}+25y''=-5+2xe^{3x}-\cos(\sqrt5 x),
\]
es: \rp\ (No calcule las constantes, sólo escriba la forma adecuada).

\vspace{0.3cm}
\noindent\textbf{c) (8\%)} Si $y(x)=c_1x^{-3}+c_2x^{-3}\ln(x)$, $x\in(0,\infty)$, es la solución general de una ecuación diferencial lineal homogénea de Cauchy-Euler de orden dos, entonces la ecuación diferencial es: \rp

\vspace{0.3cm}
\noindent\textbf{d) (10\%)} Si $y_1$ y $y_2$ son soluciones L.I. en $(0,\infty)$ de la ecuación diferencial $xy''+3y'+xe^xy=0$ y si el Wronskiano de $y_1$ y $y_2$ evaluado en 1 es $W(y_1,y_2)(1)=3$, entonces
\[
W(y_1,y_2)(5)=\rp
\]

\vfill
\rule{\linewidth}{0.4pt}\\[1pt]
{\scriptsize\textit{Transcripción digital --- MathUNAL}\hfill\textcolor{unalblue}{\textbf{1}}}

\end{document}
